\chapter{Device failures, Stop on Fail and Abort}\label{ch:errors}

\section{Two kinds of failure}

A \textbf{Device failure} is the connection itself going wrong: nothing listening on the port,
the socket reset, the serial port gone, a read that times out with no bytes. These are handled
with \texttt{onError}.

A \textbf{schedule bug} is a bad argument: a field the Device Type does not know, a missing
required field, a value of the wrong type, an \texttt{ATE} function called where it is not
allowed. That is not a Device failure and never reaches \texttt{onError}. In the code it is
isolated: the call returns \texttt{false} and the message, and the log carries an ERROR line
starting with the schedule's own file and line. ATE never lets a Python exception reach Lua.

\section{\texttt{onError}}

\texttt{ATE.open(devType, params, onError)} and \texttt{ATE.read(dev, params, onError)} take a
handler called as \texttt{onError(opResult, opArgs)}: the failure table the call would have
returned, with \texttt{err} set, and the params table that was passed in. If a later call has no
handler of its own, the one given to \texttt{open} is used.

\begin{itemize}
  \item If \texttt{onError} returns anything other than \texttt{nil}, the failed call returns
    that value and the schedule carries on.
  \item If it returns \texttt{nil}, or there is no handler, the call raises a Lua error. Unless
    the step catches it, the step records \texttt{Result.ERROR}.
  \item The error message names the function and the Device: \texttt{tcp.tslua:42: ATE.read:
    TCP 127.0.0.1:5025: timeout}.
\end{itemize}

A read that times out \emph{after} some bytes arrived is not a Device failure: it returns
\texttt{partial = true} and logs a WARNING. ATE never reconnects on its own; to recover, close
the Device and open it again with the same params, which gives a new handle.

\example{../examples/test-writer/errors.tslua}{errors.tslua}

\section{Stop on Fail and Retry}

By default a step that ends \texttt{FAIL} or \texttt{ERROR} pauses the run and the Operator
chooses \texttt{CONTINUE}, \texttt{RETRY} or \texttt{ABORT}. \texttt{:setStopOnFail(false)}
turns that off for the run, and the Operator can change it on the Run screen before starting.

A Retry runs the step's action again from the beginning, with its POLL inputs reset and its
blocking inputs cleared. Every Attempt is recorded; the step's result is its last Attempt's. A
schedule has no way to read its own Attempt number, so keep the count yourself in a local the
action closes over.

\example{../examples/test-writer/retry.tslua}{retry.tslua}

\section{Abort}

Abort is a whole-run stop. The step running when it lands records \texttt{ERROR} with the
message \texttt{Aborted by Operator}; every step that never ran records \texttt{NOTRUN}; the
run's End Reason is \texttt{ABORTED} and its result \texttt{FAIL}.

An abort cannot be caught. Once it is requested, \texttt{ATE.sleep} and \texttt{ATE.input} wake
immediately, a busy Lua loop is stopped by an instruction-count hook, and \texttt{pcall} and
\texttt{xpcall} re-raise rather than returning. ATE aims to stop within half a second and
promises within ten.

Post-test Actions still run afterwards, each with the run's result. A second abort ends only the
Action it lands in. Every Device still open is closed.

\example{../examples/test-writer/abort.tslua}{abort.tslua}

\cando{handle a Device failure without letting it end a run, and explain what Stop on Fail,
Retry and Abort each do to the record.}
